
En los siguientes apartados se encuentran los diferentes repositorios involucrados en este documento. Se puede acceder a todo el software desarrollado en este proyecto libremente.


\section{Útiles y reportes\label{SEC:UTILESYREPORTSREPO}}


En este repositorio se encuentran diversos útiles como tutoriales (para instalar \textit{Preempt-RT}, los drivers de Nvidia, los drivers de la DAQ de National Instruments, etc. Para más detalles, se puede revisar la lista de \textit{Contents} en el apartado de \textit{Tutorials and Reports}) así como implementaciones para algunos modelos neuronales en c y cpp junto con un reporte sobre la ejecución de los mismos.

Se puede acceder a través de \MYhref{https://github.com/sergiohidalgo818/tfg-utils-and-docs}{este enlace} o de este código QR (Figura \ref{FIG:UTILSREPORTS}).

\begin{figure}[QR repositorio de útiles y reportes]{FIG:UTILSREPORTS}{Código QR al repositorio de útiles y reportes (\MYhref{https://github.com/sergiohidalgo818/tfg-utils-and-docs}{link}).
    }
    \image{0.25\textwidth}{}{qr-utils-reports}
\end{figure}


\section{RTXI para Preempt-RT\label{SEC:RTXIPREEMPTREPO}}


En este repositorio se encuentran las dos implementaciones realizadas en las ramas del mismo. Una de ellas es la de \textit{Preempt-RT} (ramas `preempt-rt` y `preempt-rt-stable`, siendo la última en la que se permiten ciclos en las conexiones de los plúgins) y el botón para permitir los ciclos (rama `connector-cycle-button` y `preempt-rt-cyclic`, siendo esta última un merge de ambos cambios).

Se puede acceder a través de \MYhref{https://github.com/sergiohidalgo818/rtxi}{este enlace} o de este código QR (Figura \ref{FIG:RTXI}).

\begin{figure}[QR repositorio de RTXI]{FIG:RTXI}{Código QR al repositorio del fork de \textit{RTXI} (\MYhref{https://github.com/sergiohidalgo818/rtxi}{link}).
    }
    \image{0.25\textwidth}{}{qr-rtxi}
\end{figure}


\section{Datos y gráficas\label{SEC:DATOSREPO}}


En este repositorio se encuentran todos los datos usados para las generar las gráficas expuestas (aunque también se encuentran otras gráficas que no han sido incluidas en este documento pero que pueden ser de interés). También se encuentran los scripts usados para generar dichas gráficas, así como otros recursos relacionados con la obtención de dichos datos, como los valores usados para obtenerlos y en el caso de los circuitos de \textit{RTXI}, esquemas de como conectar los plugins. \textbf{Importante:} Este repositorio usa \MYhref{https://git-lfs.com/}{Git LFS}.

Se puede acceder a través de \MYhref{https://github.com/sergiohidalgo818/tfg-data}{este enlace} o de este código QR (Figura \ref{FIG:DATA}).

\begin{figure}[QR repositorio de datos]{FIG:DATA}{Código QR al repositorio de datos (\MYhref{https://github.com/sergiohidalgo818/tfg-data}{link}).
    }
    \image{0.25\textwidth}{}{qr-data}
\end{figure}


\section{RTHybrid para RTXI v3\label{SEC:RTHPRTXI3REPO}}


En este repositorio es un fork del original, en este se encuentran las implementaciones para \textit{RTXI V3} de los plugins de \textit{RTHybrid}. Los cambios se encuentran en la rama `rthybrid-for-rtxi3` (tanto los links como el QR redirigen ya a la rama).

Se puede acceder a través de \MYhref{https://github.com/sergiohidalgo818/rthybrid-for-rtxi/tree/rthybrid-for-rtxi3}{este enlace} o de este código QR (Figura \ref{FIG:RTHYBRID4RTXI}).

\begin{figure}[QR repositorio de los plugins de RTHybrid para RTXI V3]{FIG:RTHYBRID4RTXI}{Código QR al repositorio del fork de los plugins de \textit{RTHybrid} (\MYhref{https://github.com/sergiohidalgo818/rthybrid-for-rtxi/tree/rthybrid-for-rtxi3}{link}).
    }
    \image{0.25\textwidth}{}{qr-rthybrid-for-rtxi3}
\end{figure}


\section{Pico Neuron\label{SEC:PICONEURONREPO}}


En este repositorio se encuentra el programa realizado para ejecutar un modelo neuronal en una Raspberry Pico W y una Raspberry Pico 2.

Se puede acceder a través de \MYhref{https://github.com/sergiohidalgo818/pico-neuron}{este enlace} o de este código QR (Figura \ref{FIG:PICONEURON}).

\begin{figure}[QR repositorio de pico-neuron]{FIG:PICONEURON}{Código QR al repositorio de \textit{pico-neuron} (\MYhref{https://github.com/sergiohidalgo818/pico-neuron}{link}).
    }
    \image{0.25\textwidth}{}{qr-pico-neuron}
\end{figure}


Para medir su salida se ha creado también un programa en c con un \textit{realtime thread}.

Se puede acceder al repositorio \MYhref{https://github.com/sergiohidalgo818/pico-neuron-serial}{este enlace} o de este código QR (Figura \ref{FIG:PICONEURONSERIAL}).

\begin{figure}[QR repositorio de pico-neuron-serial]{FIG:PICONEURONSERIAL}{Código QR al repositorio de \textit{pico-neuron-serial} (\MYhref{https://github.com/sergiohidalgo818/pico-neuron-serial}{link}).
    }
    \image{0.25\textwidth}{}{qr-pico-neuron-serial}
\end{figure}


